\scanpage{39}
% archive foliation: 38
\margin{\figure{\begin{tikzpicture}\draw (0,0)--(.15,-2.8);\draw (.15,-4.6)--(.18,-5.6);\end{tikzpicture}}}
\uline{Proposition} $X_i$ ($1\leq i\leq n$) de type fini sur $Y$ de type fini sur $S$ noeth, $F_i$ coh. sur $X_i$, $N$ un entier. Alors il existe un \deleted{entier} ouvert non vide $U$ de $S$, tel que, \deleted{\unsure{désignant par}} \inserted{posant}
\[
\deleted{\unsure{U'}}\inserted{S'}=U_{\text{réd}},\quad X'_i=X\times_S\deleted{\unsure{U'}}\inserted{S'},\quad Y'=Y\times_S\deleted{\unsure{U'}}\inserted{S'},\quad F'_i=F\otimes_S S',
\]
on ait ceci : pour tout changement de base $T\longrightarrow S'$, le hom canonique
\[
\operatorname{Tor}_p^{\underline O_{S'}}(F'_1,-,F'_n)\otimes_{S'}T\longrightarrow\operatorname{Tor}_p^{\underline O_T}(F'_{1T},-,F'_{nT})
\]
(où $F'_{\alpha T}=F'_\alpha\otimes_{S'}T$) est un \uline{isom} pour $p\leq N$.

On utilise pour ceci le

\margin{\figure{\begin{tikzpicture}\draw (0,0)--(.08,-2.5);\draw (.14,-2.7)--(.2,-3.9);\draw (.25,-4.8)--(.3,-5.2);\end{tikzpicture}}}
\uline{Lemme} \inserted{Avec les notations précédentes mais sans condition de finitude,} Supposons que $Y$ soit \uline{plat} sur $S$, que les $F_\alpha$ soient plats sur $S$, et que les $\operatorname{Tor}_i^{\underline O_{\unsure{S|Y}}}(F_1,-,F_n)$ pour $i\leq N$ soient plats sur $S$. Alors pour tt changement de base $T-S$, l'hom canonique
\[
\operatorname{Tor}_i^{\underline O_S}(F_1,-,F_n)\times_S T-\operatorname{Tor}_i^{\underline O_T}(F_{1(T)},-,F_{n(T)})
\]
est un isom pour $i\leq N$.

\margin{\unsure{Peut-on faire plus ?}}
